\section*{ОБОЗНАЧЕНИЯ И СОКРАЩЕНИЯ}

API -- программный интерфейс приложения (Application Programming Interface).

C\# -- язык программирования, используемый для реализации логики приложения.

JSON (JavaScript Object Notation) -- текстовый формат обмена и хранения структурированных данных.

MVVM (Model--View--ViewModel) -- архитектурный шаблон, разделяющий модель данных, пользовательский интерфейс и модель представления.

QR-код -- двумерный штрихкод, используемый для быстрого определения точки на карте.

UI (User Interface) -- пользовательский интерфейс.

UX (User Experience) -- пользовательский опыт.

XAML (Extensible Application Markup Language) -- язык разметки пользовательского интерфейса в приложениях .NET MAUI.

Граф -- структура данных, состоящая из узлов и рёбер, используемая для представления помещений, переходов и маршрутов.

ИС -- информационная система.

ПО -- программное обеспечение.

ТЗ -- техническое задание.

ТП -- технический проект.

ЧС -- чрезвычайная ситуация.

SVG (Scalable Vector Graphics) -- формат векторной графики, используемый для представления планов этажей.

UML (Unified Modelling Language) -- язык графического описания для объектного моделирования в области разработки программного обеспечения.
